% appP_checks.tex -- generated by writing/tools/blueprint_gen.py 1.0 (level 2) on 2026-10-14 from sources 20021204b9; do not edit
\providecommand{\bpNum}[1]{\hyperref[#1]{\ref*{#1}}}%
\providecommand{\bpNumA}[1]{\hyperref[#1]{\textcolor{orange!65!black}{\ref*{#1}}}}%
\providecommand{\bpPage}[1]{\pageref*{#1}}%
\providecommand{\bpLab}[1]{{\ttfamily #1}}%
\providecommand{\bpBrk}{\allowbreak}%
\providecommand{\bpT}[2]{\facttagbox[#1]{\scriptsize #2}}%
\providecommand{\bpU}[1]{{\color{gray!70!black}#1}}%

\section{一致性检查}\label{sec:appP-checks}

\begingroup\scriptsize\setlength{\tabcolsep}{2.5pt}\renewcommand{\arraystretch}{1.06}
\begin{longtable}{@{}>{\raggedright\arraybackslash}p{\dimexpr(\textwidth-6\tabcolsep-0mm)*72/1000\relax}>{\raggedright\arraybackslash}p{\dimexpr(\textwidth-6\tabcolsep-0mm)*257/1000\relax}>{\raggedright\arraybackslash}p{\dimexpr(\textwidth-6\tabcolsep-0mm)*106/1000\relax}>{\raggedright\arraybackslash}p{\dimexpr(\textwidth-6\tabcolsep-0mm)*564/1000\relax}@{}}
\caption{一致性检查的结果（生成程序 blueprint\_gen.py 版本 1.0，日期 2026-10-14，源文件修订 \texttt{20021204b9}，压缩级别 2；E 为错误，W 为警告，I 为信息）}\label{tab:appP-checks}\\
\toprule
\textbf{检查} & \textbf{内容} & \textbf{结果} & \textbf{前三条发现} \\
\midrule
\endfirsthead
\caption[]{一致性检查的结果（续）}\\
\toprule
\textbf{检查} & \textbf{内容} & \textbf{结果} & \textbf{前三条发现} \\
\midrule
\endhead
\midrule
\multicolumn{4}{r@{}}{\footnotesize（续下页）}\\
\endfoot
\bottomrule
\endlastfoot
\mbox{C-P0} & 源文件解析（环境、标题、标签） & 通过 & — \\
\mbox{C-P1} & 依赖、前提与非微扰地位中的标签均有定义；\texttt{\string\uses} 的格式 & E\,1 & \bpLab{\detokenize{def:}\bpBrk \detokenize{intro-}\bpBrk \detokenize{tags}}：未知标签 \bpLab{\detokenize{asm:}\bpBrk \detokenize{intro}} \\
\mbox{C-P2} & 规范标签由登记的所属单元写出，且与登记表一致 & E\,178 & 规范标签 \bpLab{\detokenize{fig:}\bpBrk \detokenize{factorization-}\bpBrk \detokenize{schematic}}（Ch1）缺失；\newline 规范标签 \bpLab{\detokenize{fig:}\bpBrk \detokenize{handbag}}（Ch1）缺失；\newline 规范标签 \bpLab{\detokenize{fig:}\bpBrk \detokenize{intro-}\bpBrk \detokenize{dependency}}（Ch1）缺失 \\
\mbox{C-P3} & 无前向依赖 & 通过（I\,1） & 章到附录的依赖 180 条（允许） \\
\mbox{C-P4} & 无依赖环 & 通过 & — \\
\mbox{C-P5} & 等级传递：标签不高于前提的内容等级 & 通过 & — \\
\mbox{C-P6} & 正面陈述不依赖 [V] 陈述 & 通过 & — \\
\mbox{C-P7} & [R] 的限定语逐字一致且彼此相容 & 通过（I\,1） & 约定性假设 \bpLab{\detokenize{asm:}\bpBrk \detokenize{appN-}\bpBrk \detokenize{values}}（无标签）被使用 \\
\mbox{C-P8} & $R_{\mathrm{cond}}$ 的前提句完整 & W\,37 & \bpLab{\detokenize{lem:}\bpBrk \detokenize{method-}\bpBrk \detokenize{local-}\bpBrk \detokenize{accuracy}}：前提 \bpLab{\detokenize{lem:}\bpBrk \detokenize{appN-}\bpBrk \detokenize{error-}\bpBrk \detokenize{table}}（A）不在前提句中；\newline \bpLab{\detokenize{prop:}\bpBrk \detokenize{disall-}\bpBrk \detokenize{two-}\bpBrk \detokenize{routes}}：前提 \bpLab{\detokenize{lem:}\bpBrk \detokenize{appN-}\bpBrk \detokenize{error-}\bpBrk \detokenize{table}}（A）不在前提句中；\newline \bpLab{\detokenize{prop:}\bpBrk \detokenize{disall-}\bpBrk \detokenize{two-}\bpBrk \detokenize{routes}}：前提 \bpLab{\detokenize{lem:}\bpBrk \detokenize{method-}\bpBrk \detokenize{local-}\bpBrk \detokenize{accuracy}}（A）不在前提句中 \\
\mbox{C-P9} & 每条编号假设至少被使用一次 & 通过 & — \\
\mbox{C-P10} & 「被使用于」与计算的直接使用者一致 & 通过 & — \\
\mbox{C-P11} & 主表各行的标签存在且严格性标签一致 & 通过 & — \\
\mbox{C-P12} & 无重复标签 & 通过 & — \\
\mbox{C-P13} & 无待核实标记与已撤回的标签 & 通过 & — \\
\mbox{C-P14} & 主表“瓶颈”列包含计算的瓶颈 & W\,119 & 主表 IV.31（\bpLab{\detokenize{conj:}\bpBrk \detokenize{break-}\bpBrk \detokenize{dps-}\bpBrk \detokenize{coloured}}）：缺 13、多 0；\newline 主表 III.19（\bpLab{\detokenize{conj:}\bpBrk \detokenize{coll-}\bpBrk \detokenize{resolved-}\bpBrk \detokenize{photo}}）：缺 13、多 0；\newline 主表 III.31（\bpLab{\detokenize{conj:}\bpBrk \detokenize{dy-}\bpBrk \detokenize{coloured-}\bpBrk \detokenize{final}}）：缺 13、多 0 \\
\mbox{C-P15} & 强子层面陈述的非微扰地位行 & E\,49，W\,289 & \bpLab{\detokenize{thm:}\bpBrk \detokenize{qcd-}\bpBrk \detokenize{asymptotic-}\bpBrk \detokenize{freedom}}：强子层面陈述缺非微扰地位行；\newline \bpLab{\detokenize{thm:}\bpBrk \detokenize{ir-}\bpBrk \detokenize{leading-}\bpBrk \detokenize{regions}}：强子层面陈述缺非微扰地位行；\newline \bpLab{\detokenize{thm:}\bpBrk \detokenize{ee-}\bpBrk \detokenize{R-}\bpBrk \detokenize{finite}}：非微扰地位行含非全局标签 \bpLab{\detokenize{conj:}\bpBrk \detokenize{ee-}\bpBrk \detokenize{duality}} \\
\mbox{C-P16} & 「地位与缺口」段落存在并点名瓶颈 & W\,142 & \bpLab{\detokenize{thm:}\bpBrk \detokenize{ir-}\bpBrk \detokenize{leading-}\bpBrk \detokenize{regions}}：「地位与缺口」未点名瓶颈（\bpLab{\detokenize{prop:}\bpBrk \detokenize{ir-}\bpBrk \detokenize{pinch-}\bpBrk \detokenize{classification}} 等）；\newline \bpLab{\detokenize{thm:}\bpBrk \detokenize{sudakov-}\bpBrk \detokenize{factorization}}：「地位与缺口」未点名瓶颈（\bpLab{\detokenize{prop:}\bpBrk \detokenize{ir-}\bpBrk \detokenize{regions-}\bpBrk \detokenize{minkowski}} 等）；\newline \bpLab{\detokenize{prop:}\bpBrk \detokenize{sudakov-}\bpBrk \detokenize{kg}}：「地位与缺口」未点名瓶颈（\bpLab{\detokenize{prop:}\bpBrk \detokenize{ir-}\bpBrk \detokenize{regions-}\bpBrk \detokenize{minkowski}} 等） \\
\mbox{C-P17} & 具名缺口被点名 & W\,38 & \bpLab{\detokenize{prop:}\bpBrk \detokenize{method-}\bpBrk \detokenize{glauber-}\bpBrk \detokenize{deformation}}：未点名具名缺口 \bpLab{\detokenize{asm:}\bpBrk \detokenize{appN-}\bpBrk \detokenize{brst-}\bpBrk \detokenize{nonlocal}}；\newline \bpLab{\detokenize{conj:}\bpBrk \detokenize{coll-}\bpBrk \detokenize{resolved-}\bpBrk \detokenize{photo}}：未点名具名缺口 \bpLab{\detokenize{asm:}\bpBrk \detokenize{appN-}\bpBrk \detokenize{brst-}\bpBrk \detokenize{nonlocal}}；\newline \bpLab{\detokenize{prop:}\bpBrk \detokenize{coll-}\bpBrk \detokenize{diffractive-}\bpBrk \detokenize{photo}}：未点名具名缺口 \bpLab{\detokenize{asm:}\bpBrk \detokenize{appN-}\bpBrk \detokenize{brst-}\bpBrk \detokenize{nonlocal}} \\
\end{longtable}
\endgroup
完整的发现清单（按单元分组）见生成程序的报告文件 \texttt{appP\_report.txt}。
\FloatBarrier\label{appP:end}
